\section{Ideals Generated By A Subset}
\begin{definition}[Ideal Genereated By A Subset]
Let $A$ be a subset of a ring $R$ with identity $1 \neq 0$. Let $S = \{I_k\}$ be the collection of all left ideals of $R$ that contain $A$ (note that the set is nonempty since $A \subseteq R$ and $R$ is an ideal in itself). The intersection
\[
\bigcap_{I_k \in S} I_k,
\]
is called the \textit{left ideal generated by} $A$, and is denoted by $(A)_L$.
Alternatively, we can constructively form the set of elements that constitutes this ideal: 
\[(A)_L = \left\{\sum_{k=1}^{m}r_ka_k: r_k \in R, a_k\in A, m \in \N\right\}.\]
The definition is symmetrical for \textit{right ideal generated by a subset}.
\end{definition}
\begin{definition}[Finitely Generated Ideal]
    An ideal in a commutative ring is said to be \textit{finitely generated} if there exits $x_1, x_2, ..., x_n$ such that 
    \[I = Rx_1 + Rx_2 + \cdots + Rx_m, m \in \N.\]
    Or for all $i \in I$,
    \[i = r_1 x_1 + r_2 x_2 + \cdots + r_nx_n.\]
\end{definition}
\begin{eg}
    $R = \Z, I = 2\Z + 3\Z = (2,3) = \{2k + 3l: k,l \in \Z\}$. Hence, $2\Z + 3\Z$ is a finitely generated ideal in $\Z$.
\end{eg}
\begin{definition}[Principle Ideal]
    \textit{Principle ideal} is an ideal generated by a single element $a$. Symbolically,
    \[(a) = \{ra: r \in R\}\]
\end{definition}
\begin{eg}
    $R = \Z, I = 3\Z = (3) = \{k \cdot 3 : k \in \Z\}$. Hence, $3\Z$ is a principle ideal in $\Z$.
\end{eg}
\begin{prop}
    Let $I$ be an ideal of $R$ with identity. Then,
    \begin{enumerate}
        \item $I = R$ if and only if $I$ contains the unit element.
        \item $I = R$ if and only if $I$ contains a unit.
        \item Assume $R$ is commutative, then $R$ is a field if and only if its only ideals are $(0)$ and $(1)=R$ itself.
    \end{enumerate}
\end{prop}
\begin{propproof}
    \begin{enumerate}
        \item If $I = R$, then its obvious. Now let $I$ contains $1$, then $r \cdot 1 = r \in I$, for all $r \in R$. Therefore, $R \subseteq I$. But $I \subseteq R$ by definition. Hence, $I = R$.
        \item If $I = R$, then its obvious. Let $a$ be a unit in $R$. Then, there exits $b \in R$ such that $ba = 1$. But since $I$ is an ideal $ba = 1 \in I$. From the previous result, $I = R$.
        \item Let $R$ be a field and $I$ be a nonzero ideal of $R$ and $x \in I$. Then $x x^{-1} = 1 \in I$. Since $1 \in I$, $I = R$. Hence, only ideals of $R$ are $(0)$ and $R$ itself. Conversely, if $(0)$ and $R$ are the only ideals of $R$, let $u$ be any nonzero element of $R$. By hypothesis $(u) = R$ and so $1 \in (u)$. Thus there is some $v \in R$ such that $1 = vu$, which means $u$ is an unit. It is true for all $u \in R$. Hence, $R$ is a commutative division ring, i.e., field. 
    \end{enumerate}
\end{propproof}
